% Weekly exercise sheet for Lecture 11; also included in the synthesis sheet.
\ifdefined\ExercisesSynthesis
\else
\documentclass[11pt]{article}
\usepackage[utf8]{inputenc}
\usepackage[T1]{fontenc}
\usepackage[english]{babel}
\usepackage{lmodern,microtype}
\usepackage[margin=0.9in]{geometry}
\usepackage{amsmath,amssymb,amsthm,mathtools}
\usepackage{enumitem,hyperref}
\newcommand{\E}{\mathbb E}
\newcommand{\Qbb}{\mathbb Q}
\newcommand{\dd}{\,\mathrm d}
\newcommand{\F}{\mathcal F}
\theoremstyle{definition}
\newtheorem{exercise}{Exercise}[section]
\newlist{exparts}{enumerate}{1}
\setlist[exparts,1]{label=\textbf{(\alph*)},itemsep=0.3em,topsep=0.4em,leftmargin=2em}
\newcommand{\lecturesummary}[1]{\begin{quote}\small\textbf{Focus.} #1\end{quote}}
\title{Stochastic Calculus and Option Pricing\\Lecture 11 Exercises}
\author{}
\date{}
\begin{document}
\maketitle
\vspace{-1.2em}
\setcounter{section}{10}
\fi

\section{Week 11 --- Change of Numeraire}
\lecturesummary{The exercise tests the no-arbitrage structure of currency
quotes.}


\begin{exercise}[Three currency quotes must agree]
Let \(X^{A/B}\) mean units of currency \(A\) per unit of currency \(B\), and
define \(X^{B/C}\) and \(X^{A/C}\) similarly.
\begin{exparts}
  \item Write the spot-quote identity that rules out a triangular arbitrage.
  Explain the units on both sides.
  \item If the identity fails, describe the cycle of conversions that
  produces a positive terminal amount from one unit of currency.
  \item Explain why the arbitrage argument uses traded exchange rates
  directly; no choice of numeraire can repair inconsistent spot quotes.
\end{exparts}
\end{exercise}

\ifdefined\ExercisesSynthesis\else
\end{document}
\fi
